% !TeX root = ../main.tex
% Read docs/writing-guide.md for purpose, outline and common mistakes.
\chapter{Verification of Prototype}
\label{chap:verification_of_prototype}

% Present results first; discuss their implications and limits separately.
\section{Illustrative results}
\label{sec:results}
Table~\ref{tab:results} shows all six constructed durations.
Concept A has a mean of \qty{60}{\second}; concept B has a mean of
\qty{65}{\second}. The observations are fictional, not experimental evidence.
% l = left-aligned text, r = right-aligned number; & separates cells, \\ ends a row.
% booktabs supplies horizontal rules. Keep the whole table on one page.
\begin{table}[htbp]
\centering
\begin{tabular}{lrrrr}
\toprule
Concept & Trial 1 (s) & Trial 2 (s) & Trial 3 (s) & Mean (s)\\
\midrule
A & 59 & 60 & 61 & 60\\
B & 64 & 65 & 66 & 65\\
\bottomrule
\end{tabular}
\caption{Fictional durations for two imaginary concepts; every entry is constructed for teaching.}
\label{tab:results}
\end{table}


\section{Verification and limitations}
Using Equation~\ref{eq:mean}, concept A meets R1 in
Section~\ref{sec:requirements}, whereas concept B does not.
Only the illustrative mean-duration criterion is assessed. The tiny invented
sample cannot quantify reliability or support inference about real timers.
A real evaluation would report uncertainty, repeatability and potential bias.
